\metatitle{Real-rooted polynomials: Catalan families}

\metadescription{Real-rootedness and interlacing for Narayana polynomials, Motzkin polynomials, Dyck path statistics, and related Catalan families.}


\section[realRootedCatalan]{Sturm sequences: Catalan families}

This page collects examples of \hyperref[sturmSequence]{Sturm sequences} arising from Catalan combinatorics:
Dyck paths, ballot sequences, Motzkin paths, and related families.
Each sequence is a \hyperref[sturmSequence]{Sturm sequence}, satisfying $P_j \interl P_{j+1}$;
in particular every polynomial in the sequence is real-rooted.
For background on interlacing and the proof techniques,
see \hyperref[interlacingPolynomials]{interlacing polynomials}.
For the broader context,
see \hyperref[polynomialRootIntro]{unimodality and real-rootedness}.

\begin{example*}[Narayana polynomials]\label{ex:narayanaSturm}

Let $P_n(t)\coloneqq \sum_{\pi \in Dyck(n)} t^{peaks(\pi)}$,
where $Dyck(n)$ Dyck paths of size $n$, and $peaks(\pi)$ count the number of peaks.
Equivalently, this is the same as counting descents in binary words with $n$ 0s and $n$ 1s,
where in each prefix of the word, there are at least as many $0$s as $1$s, see \oeis{A001263}.
A closed-form formula is
\[
P_n(t) = \frac{1}{n} \sum_{k=1}^n \binom{n}{k} \binom{n}{k-1} t^k
\]
These satisfy the recursion (see \cite[Eq. (2)]{Sulanke2002}, \cite[Eq. (3.3)]{LiuWang2007})
\[
(n+1) P_n(t) = (2n-1)(1+t)P_{n-1}(t) - (n-2)(1-t)^2 P_{n-2}(t),
\]
with initial conditions  $P_1(t)=t$, $P_2(t)=t(1+t)$, which is used to prove
the interlacing property.

Real-rootedness and the interlacing property was first proved in \cite{Branden2006}.
See also \cite[Eq (31)]{Agapito2014}, where a different recursive formula (involving derivatives)
is listed:
\[
 (n+1)(n+2) P_{n+1} = t((1-t)^2 P''_{n} + 2(1+2n)(1-t)P'_{n} + (1+2n)(2+2n)P_n).
\]
Another approach is given in \cite{KostovFinkelshteinShapiro2009}.

\begin{lstlisting}

NN[1] := t;
NN[2] := t (1 + t);
NN[n_] := NN[n] = Expand[
	(
		(2 n - 1) (1 + t) NN[n - 1] -
		(n - 2) (1 - t)^2 NN[n - 2]
)/(n + 1)];
\end{lstlisting}

One can also prove real-rootedness by observing that the Narayana polynomials
are a simple transformation of Jacobi polynomials, or Gegenbauer polynomials
(which are orthogonal polynomials on an interval).
We have that the Narayana polynomials are given by the Hypergeometric series:
\[
P_n(t) = {}_2F_1(-n, -n-1;2;x)
\]
or, \texttt{HypergeometricPFQ[{-n, -n - 1}, {2}, x]}, see \cite[Sec. 12]{BarryHennessy2011}.


Yet another recurrence is
\[
(1 + n)P_n=(3nt + n - t - 1) P_{n-1} - 2 t(t-1) P'_{n-1}
\]
which can be proved using the first theorem in \cite{Dung2022x}.

\name{Shi-Mei Ma}, \name{Jean Yeh}, and \name{Yeong-Nan Yeh} give a
context-free grammar for the Narayana polynomials \cite{MaYehYeh2023x}.
Their differential-operator framework also gives standard-Young-tableau
interpretations for several Eulerian and Narayana-type polynomial families.

\name{Harold R. L. Yang} and \name{Philip B. Zhang} introduce stable
multivariate refinements of the type $A$ and type $B$ Narayana polynomials
\cite{YangZhang2024x}.  Their polynomials are generating functions for
labeled plane trees with weights recording improper edges, and real stability
is proved using the Borcea--Brändén characterization of stability-preserving
linear operators.
\name{Frédéric Chapoton} and \name{Guo-Niu Han} study the roots of the
Poupard and Kreweras polynomials \cite{ChapotonHan2020}.
\end{example*}


\begin{example*}[Type $B$ Narayana polynomials]
The type $B$ Narayana polynomials can be defined as
\[
P_n(t) = \sum_{k=0}^n \binom{n}{k}^2 t^k.
\]
Since these are produced by the Hadamard product of $(1+t)^n$ with itself, it
follows that the $P_n(t)$ are also real-rooted.
See the \hyperref[hadamardProduct]{Hadamard product subsection} for a general statement.
\end{example*}



\begin{example*}[Generalized Narayana polynomials]
A two-parameter family of Narayana-type polynomials
\[
P_{n,m}(t) = \sum_{k=0}^n \left(\binom{n}{k}\binom{m}{k} -\binom{n}{k+1}\binom{m}{k-1} \right)t^k.
\]
is considered in \cite[Thm. 1.4]{ChenYangZhang2018},
and they use a fairly general recurrence relation to obtain interlacing roots.
\end{example*}


\begin{example*}[Boolean--Narayana numbers]

A family similar to the  Narayana numbers, named the \defin{Boolean--Narayana numbers}
is introduced by \name{Miklós Bóna} in \cite{Bona2026x}.
The corresponding polynomials are proved to be real-rooted.
\[
 BoNa(n,k) = \frac{1}{n}\binom{n}{k-1}\sum_{j=0}^{k-1}2^j\binom{k-1}{j}\binom{n-k+1}{j+1}.
\]
If we define $P_n(x) \coloneqq \sum_{k} BoNa(n,k)x^k$ then these satisfy the recursion
\[
 (n+1) P_n(x) = (2n-1)(1+x) P_{n-1}(x) - (n-2)(1-6x+x^2) P_{n-2}(x).
\]
\end{example*}


\begin{example*}[NNE instances in Dyck paths with fixed number of peaks]

In \cite{WangZhang2024}, the authors consider the polynomials
\[
 W_{n,k}(t) \coloneqq [x^k] \sum_{\pi \in Dyck(n)} x^{\mathrm{peaks}(\pi)} t^{\mathrm{nne}(\pi)},
\]
where $\mathrm{nne}(\pi)$ count the number of North-North-East instances in the path.
The authors show that $\{ W_{n,k}(t) \}_{n \geq k}$ for fixed $k$ form a \hyperref[sturmSequence]{Sturm sequence},
and that $\{ W_{n,k}(t) \}_{1\leq k \leq n}$ is \hyperref[sturmSequence]{Sturm-unimodal}.

\begin{conjecture}[\name{Per Alexandersson}, 2023]
The generalization to Fuss--Catalan paths (in an $n \times nm$ rectangle)
seem to be real-rooted as well.
\end{conjecture}


The following appears in \cite{BonaDimitrovLabelleLiPappeVindasMelendezZhuang2025}.

Let $w_{n,k,m}$ be the number of Dyck paths of size $n$, with $k$
instances of $NE$ and $m$ instances of $NNE$. Then the polynomials defined via
\[
W_{n,k}(t) = \sum_{m\geq 0} w_{n,k,m}t^m = \frac{1}{k} \binom{n}{k-1} \sum_{m \geq 1}^{\min(k,n-k)}  \binom{n-k-1}{m-1}\binom{k}{m} t^m
\]
are real-rooted. The formula above is valid for $n \gt k$, and for $n \leq k$ the polynomials are constant.
The authors conjectured that for fixed $k$, $\{W_{n,k}\}_{n\geq k}$ form a \hyperref[sturmSequence]{Sturm sequence},
and that for fixed $n$, $\{W_{n,k}\}_{1 \leq k\leq n}$ is \hyperref[sturmSequence]{Sturm-unimodal}.
These conjectures were later proved in \cite{WangZhang2024}.
The authors there finds a two-term recurrence, expression $W_{n,k}(t)$ in terms of $W_{n-1,k}(t)$ and $W_{n,k}(t)$.
Alternatively, $W_{n,k}(t)$ can be expressed in terms of $W_{n,k-1}(t)$ and its derivative.
\end{example*}


\begin{example*}[Motzkin polynomials]

The non-crossing matching polynomial of the complete
graph $K_n$ is the Motzkin polynomial, see \oeis{A055151}.
These polynomials satisfy the recursion
\[
 (n+2)M_n(x) = (2n+1)M_{n-1}(x) - (n-1)(1-4x)M_{n-2}(x), \text{ with } M_1(x)=1, \;
 M_2(x)=x+1.
\]
By using the methods in \cite[Section 3]{Wagner1992},
it is now easy to use induction to show that the roots of $M_{n-1}(x)$ interlaces
those of $M_{n}(x)$. In particular, these polynomials are real-rooted.

% M[1] = 1;
% M[2] = 1 + x;
% M[n_] :=
%   Together[
%     1/(n + 2) ((2 n + 1) M[n - 1] - (n - 1) (1 - 4 x) M[n - 2])] //
%    Expand;

\end{example*}
